% =====================================================================
\section{Giới thiệu}
\label{sec:intro}

CLIP~\cite{radford2021clip} đưa ảnh và câu vào cùng một không gian vector, nên truy vấn văn bản--ảnh chỉ còn là tìm ảnh
có cosine lớn nhất với câu. Cách so khớp này có hai giới hạn. Thứ nhất, nó không dùng quan hệ có cấu trúc giữa các vật
thể trong ảnh, hay giữa các thực thể trong một đồ thị tri thức. Thứ hai, nó không cho biết vì sao một ảnh được xếp lên
đầu. Đề bài yêu cầu một hệ thống trả lời ba câu hỏi nối tiếp nhau:
(1) ảnh nào phù hợp nhất với truy vấn;
(2) tri thức quan hệ giúp truy vấn tốt hơn ở đâu, nhất là với truy vấn gián tiếp;
(3) vì sao mô hình trả về kết quả này.
Nhóm làm nhánh A của đề (văn bản tìm ảnh) và thực hiện đủ ba cấp độ.

\textbf{Cấp độ~1} (Visual Genome $\cap$ COCO). CLIP đóng băng chọn 50 ứng viên. Một bộ mã hoá GATv2~\cite{brody2022gatv2}
chạy trong scene graph của từng ảnh và trong đồ thị tách từ câu, rồi chấm điểm khớp ở hai mức: vật thể và bộ ba (chủ thể,
quan hệ, đối tượng). Điểm cuối kết hợp cosine CLIP với điểm đồ thị theo dạng
$\alpha\cdot\text{CLIP} + (1-\alpha)\cdot\text{đồ thị}$ của đề. Nhóm so với CLIP thuần trên cùng cách chia. Các ablation
gồm bỏ cạnh, nối lại cạnh ngẫu nhiên và thay đồ thị của ảnh khác. Phần cuối phân tích 5 truy vấn thành công và 5 truy vấn
thất bại.

\textbf{Cấp độ~2} (thêm MKG-W). Nhóm tái lập NativE~\cite{zhang2024native}, một mô hình hoàn thiện đồ thị tri thức đa
phương thức năm 2024, và kiểm tra lại con số link prediction của bài báo. Sau đó nhóm đưa ba loại tín hiệu đồ thị vào cùng
pipeline xếp hạng lại: số cạnh tri thức, embedding của NativE, và điểm bộ ba của chính NativE. Nhóm xây thêm bộ 25 truy vấn
tự viết, tách nhóm trực tiếp và gián tiếp.

\textbf{Cấp độ~3}. Nhóm đo fidelity, validity và sparsity của lời giải thích trên toàn bộ tập validation, kèm bằng chứng
trên từng truy vấn. Nhóm cũng xây một đồ thị tri thức đa phương thức (MMKG) tiếng Việt nhỏ, chạy phương pháp của cấp độ~2
trên đó và phân tích khả năng chuyển giao.

Mục~\ref{sec:data} mô tả dữ liệu và các quy tắc chống rò rỉ. Mục~\ref{sec:method} mô tả mô hình và cách lời giải thích được sinh ra. Các
mục~\ref{sec:l1}--\ref{sec:vnmmkg} trình bày kết quả theo từng cấp độ. Mục~\ref{sec:conclusion} nêu hạn chế. Mục~\ref{sec:contrib}
ghi phân công và phần đã dùng công cụ AI. Mọi con số đều lấy từ một file kết quả trong thư mục \file{experiments/} của
repo; đường dẫn ghi ngay cạnh bảng.
